\documentclass[review]{elsarticle}
\usepackage{amsmath,amsthm,amssymb}
\usepackage{mathtext}
\usepackage[T1,T2A]{fontenc}
\usepackage[T2A]{fontenc}
\usepackage[utf8]{inputenc}
\usepackage[english, russian]{babel}
\usepackage{lineno,hyperref}
\modulolinenumbers[5]

\journal{Journal of \LaTeX\ Templates}

%%%%%%%%%%%%%%%%%%%%%%%
%% Elsevier bibliography styles
%%%%%%%%%%%%%%%%%%%%%%%
%% To change the style, put a % in front of the second line of the current style and
%% remove the % from the second line of the style you would like to use.
%%%%%%%%%%%%%%%%%%%%%%%

%% Numbered
%\bibliographystyle{model1-num-names}

%% Numbered without titles
%\bibliographystyle{model1a-num-names}

%% Harvard
%\bibliographystyle{model2-names.bst}\biboptions{authoryear}

%% Vancouver numbered
%\usepackage{numcompress}\bibliographystyle{model3-num-names}

%% Vancouver name/year
%\usepackage{numcompress}\bibliographystyle{model4-names}\biboptions{authoryear}

%% APA style
%\bibliographystyle{model5-names}\biboptions{authoryear}

%% AMA style
%\usepackage{numcompress}\bibliographystyle{model6-num-names}

%% `Elsevier LaTeX' style
\bibliographystyle{elsarticle-num}
%%%%%%%%%%%%%%%%%%%%%%%

\begin{document}

\begin{frontmatter}

\title{Elsevier \LaTeX\ template\tnoteref{mytitlenote}}
\tnotetext[mytitlenote]{Fully documented templates are available in the elsarticle package on \href{http://www.ctan.org/tex-archive/macros/latex/contrib/elsarticle}{CTAN}.}

%% Group authors per affiliation:
\author{Elsevier\fnref{myfootnote}}
\address{Radarweg 29, Amsterdam}
\fntext[myfootnote]{Since 1880.}

%% or include affiliations in footnotes:
\author[mymainaddress,mysecondaryaddress]{Elsevier Inc}
\ead[url]{www.elsevier.com}

\author[mysecondaryaddress]{Global Customer Service\corref{mycorrespondingauthor}}
\cortext[mycorrespondingauthor]{Corresponding author}
\ead{support@elsevier.com}

\address[mymainaddress]{1600 John F Kennedy Boulevard, Philadelphia}
\address[mysecondaryaddress]{360 Park Avenue South, New York}

\begin{abstract}
Привет!
\end{abstract}

\begin{keyword}
\texttt{elsarticle.cls}\sep \LaTeX\sep Elsevier \sep template
\MSC[2010] 00-01\sep  99-00
\end{keyword}

\end{frontmatter}

\linenumbers

\section{The Elsevier article class}




\section{Вступление}


В данной работе ставится \textbf{цель} проанализировать существующие методы, методологии и принципы проектирования, и на основе этого анализа синтезировать новую модель для разработки интеллектуальной системы для проектирования технического, лингвистического, программного обеспечения в области автоматизации систем управления. 

\textbf{Объектом исследования} - процесс проектирования программного обеспечения для автоматизации технологических систем :
\begin{enumerate}
  \item по типу транзакционной обработки данных,
  \item по типу предоставления доступа.
\end{enumerate}

\textbf{Предмет исследования} :
\begin{enumerate}
\item методы  создания, эффективной организации и ведения специализированного информационного и программного обеспечения АСУТП,
\item средства и методы проектирования и разработки технического, математического, лингвистического и других видов обеспечения АСУ.
\end{enumerate}


Изучаемая система может быть описана в следующем виде:
\begin{equation}
    \label{eq:equation29}
    S =  (X, T, R, Z ) ,
\end{equation}
\begin{itemize}
    \item где $ X $  — множество переменных, 
    \item $T$ — множество параметров, 
    \item $R$ — отношения на множества $X$ и $T$, 
    \item $Z$ — цель исследований т.е. решение проблемы проектирования устойчивой архитектуры программного обеспечения автоматизированной системы управления.
\end{itemize}

Перейдем к математической модели системы.
Перед вами слайд с обобщенными алгоритмическими шагами для
построения системы.


    \begin{enumerate}
        \item \textbf{Формирование базы знаний}
              \begin{itemize}
                  \item Синтез обучающей выборки
                  \item Ранжирование параметров
              \end{itemize}
        \item \textbf{Синтез деревьев и классификация}
              \begin{itemize}
                  \item Расчет гипергеометрического распределения
                  \item Синтез деревьев решений
                  \item Поиск закономерностей критерием максимума взвешенной информативности
                  \item Взвешенное голосование и классификация
              \end{itemize}
        \item \textbf{Формирование рекомендаций}
              \begin{itemize}
                  \item Формирование рекомендаций
                  \item Таблица противоречий методом ТРИЗ
              \end{itemize}
    \end{enumerate}


Далее на слайде представлена блок схема модели работы системы. Рассмотрим внимательнее данные алгоритмические блоки.

Пусть множество типов информационных систем, описываемых пользователем, как 
\begin{equation}
     U=\{u_1,u_2,...,u_n\},
\end{equation}
   множество объектов, т.е. предлагаемых архитектур 
    \begin{equation}
    P=\{p_1,p_2,...,p_n\},
    \end{equation}
    \begin{itemize}
\item     весовая матрица  $R=r_{i,j}$ размера $n$ $x$ $m$,  $i_{1...n},j_{1...m}$. 
 \item    $N$ - желаемое количество рекомендаций, которые нужно получить от системы. 
 \end{itemize}



    Требуется найти: для описываемого набора типов информационных систем $u$, найти $N $- мерный вектор 
    \begin{equation}
    \langle p_{i1},p_{i2},...,p_{iN} \rangle, 
    \end{equation}
        \begin{itemize}
        \item где архитектура $p_{i_k,k}\in N$  подходит по параметрам, т.е. в матрице описания транзакций есть пустое место $r_i,i_k$,
        \item где эти архитектуры наиболее точно соответствуют базе конюнкционных закономерностей, то есть классам, где $r_i,i_k$ является наибольшим. 
    \end{itemize}
     



Входные данные поступают в виде набора линейно-зависимых, линейно-независимых классов параметров (каналов наблюдения): вычислительные ресурсы (серверное оборудование, сеть, оборудование сопряженное, коммутация, техничесие средства ЭВМ и т .д. 

Взаимодействия между параметрами измеряется согласно следующей формуле по алгоритму квадратичной регрессии:

\begin{equation}
    \label{eq:equation_FM}
    \hat{r}_{pc} = w_{0} + \sum_{j=1}^{n} w_{j}x{pcj} + \sum_{j=1}^{n}\sum_{k=1}^{n} w_{jk}x_{pcj}x_{pck}
\end{equation}

Веса определяются низкоранговым матричным разложением:

\begin{equation}
    \label{eq:equation_wight}
    w_{jk} = x_{j}^{T}v_{k}, v_{i} \in \mathbb{R}
\end{equation}



Степень принадлежности определяется на основе рангов принадлежности и определяется на основе соотношений следующей системы:



Линейная модель ранжирования:
\begin{equation}
    \label{eq:equation_linear}
    a(x,w) = \langle x, w \rangle, x \mapsto (f_{1}(x), f_{2}(x), ..., f_{n}(x)) \in \mathbb{R},
\end{equation}

где $x \mapsto (f_{1}(x), f_{2}(x), ..., f_{n}(x)) \in \mathbb{R}$ - вектор признаков объекта $x$.

Обучающая выборка $(x_{i}, y_{i})^l_{i=1}$, где $y_{i} \in Y = {1<2<...<R}$. Функция ранжирования с параметрами w и порогами $b_0 = -\infty$, $b_{1} \leq ... \leq b_{R-1}, b_{R} = +\infty$:

\begin{equation}
    \label{eq:equation_rank_function}
    a(x,w,b) = y, b_{y-1} < g(x,w) \leq b_{y}
\end{equation}

Ранговая модель, как частный случай линейной модели $g(x,w) = \langle w, x\rangle $, сумма рассчитывается по двум порогам с функцией потерь вида:

\begin{equation}
    \label{eq:equation_loss_function3}
    \sum_{i=1}^{l}(1 - \langle w_{i}, x \rangle + b_{y_{i} - 1})_{+} + (1 + \langle x_{i}, w \rangle - b_{y_i})_{+} + \frac{1}{2C}\| w \|^2  \rightarrow min_{\substack{w,b}} 
\end{equation}

Задача квадратичного программирования, к которому можно привести выражение примет вид:

\begin{equation}
\label{eq:equation_loss_function4}
\begin{alignedat}{2}
\left\{
\begin{array}{rl}
    \frac{1}{2}\|w\|^2 + C \sum_{i=1}^{l}[y_i \neq 1]  \xi^{*}_{i} + [y_{i} \neq R]\xi_{i} \rightarrow min_{\substack{w,b,\xi}} \\
    \langle x_{i}, w \rangle \geq b_{y_{i} - 1} + 1 - \xi^{*}_{i}, \xi^{*}_{i} \geq 0, y_{i} \neq 1 ; \\
    \langle x_{i}, w \langle \geq b_{y_{i}} - 1 + \xi_{i}, \xi_{i} \geq 0, y_{i} \neq R ;
\end{array}
\right.
\end{alignedat}
\end{equation}

Двойственная задача $(\lambda^{*}_{i} = 0$ при $y_{i} = 1, \lambda_{i} = 0$ при $y_{i} = R)$


Модель ранжирования после решения двойственной задачи:

\begin{equation}
    \label{eq:rank_model1}
     \langle w, x \rangle \ = \sum_{i=1}^{l}(\lambda_{i}^{*} - \lambda_{i})K(x, x_{i}) 
\end{equation}

В такой задаче, функционал является выпуклым, как следствие имеет единственное решение. Множитель $K(x, x_{i})$ отображает возможные нелинейные обобщения $K(x, x_{i}')$. Решение будет разреженным и зависить только от опорных векторов:


где a (x,w) - параметрическая модель ранжирования.
Важным пунктом является, что параметрическая модель ранжирования влияет на деление ветвей внутри дерева. На рисунке представлен Пример работы алгоритма синтеза деревьев представлен на этом слайде.

 
Для синтеза деревьев важным критерием выступают коэффициенты ранжирования, т.к. они прямым образом влияют на ветвеление деревьев.

Пусть $X$ - вероятностное пространство , выборка $X^{l}$ - простая, то есть случайная, независимая, одинаково распределенная, $y^{*}(x), \phi(x)$ - случайные величины.
Допустим, что гипотеза о независимости событий справедлива:

\begin{equation}
    \label{eq:equation35}
    \{x: y^*(x) = c\} , \{x: \phi(x) = 1\}
\end{equation}

Тогда вероятность реализации пары $(p, n)$ возможно выразить через следующее гиперпараметрическое распределение:

\begin{equation}
    \label{eq:equation36}
     h\Bigg( 
    \begin{matrix}
    p & n \\
    P & N
    \end{matrix}
     \Bigg)
    =
    \frac{C_P^p C_N^n}{C^{p+n}_{P+N}}, 0 \leq p \leq P, 0 \leq n \leq N  ,
\end{equation}

где

\begin{equation}
    \label{eq:equation37}
    C^{k}_m = \frac{m!}{k!(m-k)!}, 0 \leq k \leq m.
\end{equation}

Таким образом, если событие с малой вероятностью реализовалось, то, скорее всего, оно не случайно. Неслучайное событие это и есть закономерность.

Основным постулатом формирования закономерности является следующее утверждение: \textit{если вероятность реализации пары совместных событий, подчиняется гипергеометрическому распределению и она мала, и тем не менее пара событий реализовалась, то гипотеза о независимости отвергается.}\\

Перейдем к результатам оценки качества работы модели.
На слайде представлен анализ параметров работы модели, первичный анализ параметров предалгается проводить согласно применению матрицы корреляций между параметрами внутри одного класса по обучающей выборке. Таким образом, можно применять методологию ТРИЗ в части выявления противоречий между параметрами.  
В диссертационной работе в Главе 4 вы можете подробно ознакомиться с результатами ранжирования для всех параметров.

На слайде 17 с синтезом деревьев расчеты статистического расчета критерия принятия или отвержения гипотезы о существовании закономерности между параметрами и отнесения объекта к некоторому классу.

На слайде 18 изображен пример синтеза дерева. Обратите, пожалуйста, внимание, что "ветвеление или разделение" дерева происходит строго по ранжированной оценке параметров. В зависимости от значимости весового коэффициента для класса, принимается решение разделить или не разделить ветвь.

На слайде 19 представлен представлены результаты ранжирования и взаимодействия между параметрами. Рабочая характеристика приемника - метрика, специальный график, позволяющий оценить качество классификации. Т.к. в данной работе велась классификация 24 классов, то суммарно было произведено 276 попарных сочетаний для сравнения рабочей характеристики приемника. В Приложении в диссертационной работы, вы сможете найти результаты каждого попарного сравнения.




\end{document}

\end{document}